\noindent\textbf{(a) Resultado.} \[u(x,y)=\sum_{n=1}^{\infty}b_n\frac{\sinh(n\pi y/a)}{\sinh(n\pi b/a)}\sin\frac{n\pi x}{a},\qquad
b_n=\frac2a\int_0^a g(s)\sin\frac{n\pi s}{a}ds.\]

\noindent\textit{Desarrollo.} Separar \(u=XY\) y usar los laterales nulos da \(X_n=\sin(n\pi x/a)\),
\(Y_n''-(n\pi/a)^2Y_n=0\). La base \(y=0\) exige \(Y_n\) proporcional a \(\sinh(n\pi y/a)\).
Dividimos por \(\sinh(n\pi b/a)\) para que el factor vertical valga \(1\) en \(y=b\), y proyectamos g en senos.
Cada término es armónico y satisface los tres bordes nulos. En todo subrectángulo con \(y<b\),
los factores verticales amortiguan exponencialmente y permiten verificar la ecuación por derivación.
La regularidad supuesta y \(g(0)=g(a)=0\) permiten recuperar g continuamente en la frontera superior.
El principio del máximo aplicado a la diferencia de dos soluciones da unicidad en el rectángulo.
No se intercambian los papeles de las coordenadas: el seno impone los laterales y la hiperbólica
lleva los datos desde \(y=b\) hacia el interior.

\medskip

\noindent\textbf{(b) Resultado.} \[u(x,y)=\frac{4a}{\pi^2}\sum_{n=1}^{\infty}\frac{\sin(n\pi/2)}{n^2}
\frac{\sinh(n\pi y/a)}{\sinh(n\pi b/a)}\sin\frac{n\pi x}{a}.\]

\noindent\textit{Desarrollo.} Partimos la integral de \(b_n\) en \(a/2\) y usamos
\(\int x\sin(kx)dx=-x\cos(kx)/k+\sin(kx)/k^2\), con \(k=n\pi/a\).
Los términos de coseno se cancelan y
\[b_n=\frac{4a}{n^2\pi^2}\sin\frac{n\pi}{2}.\]
Solo aparecen índices impares con signos alternos. La suma en \(y=b\) es exactamente
la expansión seno del triángulo; en \(y=0\) o \(x=0,a\) vale cero. Como \(0\le g\le a/2\),
la solución completa satisface \(0\le u\le a/2\) por el principio del máximo. El crecimiento aparente
de \(\sinh\) en \(y\) no produce divergencia: el cociente es \(\le1\) para \(0\le y\le b\).
